\documentclass[11pt]{article}
\usepackage[T1]{fontenc}
\usepackage{lmodern,microtype,amsmath,amssymb,amsthm,geometry,hyperref}
\geometry{margin=1in}
\hypersetup{colorlinks=true,linkcolor=blue,citecolor=blue,urlcolor=blue,
pdftitle={Pure-loss bosonic dynamic capacities: EPnI consequences and consumed-key secrecy accounting},pdfauthor={Alec Kriebel}}
\newtheorem{theorem}{Theorem}
\newtheorem{lemma}[theorem]{Lemma}
\newtheorem{proposition}[theorem]{Proposition}
\newcommand{\Tr}{\operatorname{Tr}}
\newcommand{\Loss}{\mathcal L}
\title{Pure-loss bosonic dynamic capacities:\\EPnI consequences and consumed-key secrecy accounting}
\author{Alec Kriebel\\\small ORCID: \href{https://orcid.org/0009-0001-9320-500X}{0009-0001-9320-500X}}
\date{October 6, 2026}
\setlength{\emergencystretch}{2em}
\begin{document}
\maketitle
\begin{abstract}
The finite-energy multimode entropy photon-number inequality proved in
the OpenAI mathematics release supplies the sharp entropy premise in the
previously conditional pure-loss bosonic dynamic-capacity converses.
We give its precise classical, quantum and entanglement consequence under an
average photon constraint, including arbitrary entanglement across uses.
For public, private and secret-key resources, an operational correction is
necessary: the printed catalytic secrecy condition of Wilde--Hsieh
protects consumed key jointly with generated secrets and excludes the
one-time pad. A zero-energy converse disproves the displayed region
under that literal condition. We establish the formula with an explicit
generated-secret secrecy criterion, corrected converse accounting, and
direct key-assisted coding. The formulas and original conditional
reductions are inherited; this is a consequence and operational audit,
not an independent proof of EPnI. We rely on the cited unrefereed EPnI
proof. Its analytical arguments and formal-source semantics were
audited, but its Lean build was not reproduced.
\end{abstract}

\section{Scope and attribution}
Let $\Loss_\eta$ be the one-mode pure-loss bosonic channel with an
independent vacuum environment and $\eta\in[1/2,1]$.
Its complement is $\Loss_{1-\eta}$ up to phase. For $n$ uses the
constraint is
\begin{equation}\label{energy}
 \Tr G_n\bar\rho_n\le nN,\qquad
 G_n=\sum_{j=1}^n a_j^\dagger a_j,\qquad 0\le N<\infty,
\end{equation}
averaged over classical messages, shared key and encoder randomness;
for quantum messages the logical input is maximally mixed, equivalently
the transmitted half of a maximally entangled reference pair. Arbitrary
internal multimode entanglement is allowed. No peak or per-codeword
energy constraint is asserted; noiseless side resources are charged by
rate rather than optical energy.

All logarithms and entropies are base two. Put
$g(u)=(u+1)\log_2(u+1)-u\log_2u$, $g(0)=0$.
Dividing the natural-log EPnI entropy and its $g$ by $\ln2$ leaves
the photon-number inverse unchanged. We do not mix the occasional
natural-log notation of \cite{WHG} with its bit-rate formulas.

Positive rates generate resources and negative rates consume them.
Communication resources are forward, without free feedback.
We use the original finite gross generation/consumption rates and
finite resource registers of exponential dimension, with vanishing
errors, rather than unbounded superlinear catalysts. The
$(C,Q,E)$ unit conversions are teleportation $(-2,1,-1)$,
superdense coding $(2,-1,-1)$ and entanglement distribution
$(0,-1,1)$. For $(R,P,S)$ they are one-time pad $(-1,1,-1)$,
private-to-public $(1,-1,0)$ and key distribution $(0,-1,1)$.

Wilde--Hayden--Guha \cite{WHG} gave both formulas, Gaussian ensembles
and conditional converses; De Palma \cite[Theorem 11 and Corollary 12]{DP}
gave a convex-entropy-bound formulation. The entropy breakthrough
belongs to OpenAI \cite{OAI}, at the pinned revision
\begin{center}\small\texttt{adc7f1241b42e322a6451854ab7e4b4c146bf78a}.\end{center}
The formulas already appear in arXiv:1105.0119v1 (April 30, 2011),
and the unrestricted EPnI-conditional converses in v2 (May 19, 2011).
OpenAI's public release is dated October 6, 2026; the September label
in the manuscript directory is not asserted as a public-priority date.
Our contribution is an explicit consequence and operational correction.
We claim neither an independent solution of EPnI nor first discovery of
these formulas or known coding machinery.

\section{Entropy transfer to arbitrary input ensembles}
\begin{theorem}[External EPnI consequence]\label{sharp}
For finite $n\ge1$, any finite-energy $n$-mode state $\rho$ and
$t\in[0,1]$,
\[
 \frac{H(\Loss_t^{\otimes n}(\rho))}{n}
 \ge f_t(H(\rho)/n),\qquad f_t(s)=g(tg^{-1}(s)).
\]
\end{theorem}
\begin{proof}
Substitute the independent $n$-mode vacuum into the second port of
\cite[Theorem 1.1]{OAI}, apply increasing $g$, and change entropy
units. No product, Gaussian, full-rank or finite-support condition
is added to the first input.
\end{proof}
The dependency supplement audits the upstream interpolation,
regularized minimum, Hessian, weak convolution and stationarity.
Its actual Lean declaration
\nolinkurl{OAI.EntropyPhotonNumber.entropy_photon_number_inequality}
uses full Fock states and explicit channel/entropy definitions.
The comparator's intentional \texttt{sorry} is not proof evidence.
Dependency setup exhausted disk space before kernel compilation or an
axiom report completed. No reproduced machine verification is claimed
for that input or these capacities.

\begin{lemma}\label{scalar}
$f_t$ is increasing and convex for $0\le t\le1$, and increasing
and concave for $t\ge1$.
\end{lemma}
\begin{proof}
The cases $t=0,1$ are immediate. For $u>0$,
$f_t'(g(u))=t g'(tu)/g'(u)$. The sign of its derivative is the
sign of $(tu+1)g'(tu)-(u+1)g'(u)$.
The function $(u+1)\ln(1+1/u)$ is decreasing, since its
derivative is $\ln(1+1/u)-1/u<0$.
Continuity extends the conclusions to zero.
\end{proof}
For an arbitrary ensemble $\{p_x,\rho_x\}$ satisfying \eqref{energy}
write $B_x=\Loss_\eta^{\otimes n}(\rho_x)$,
$E_x=\Loss_{1-\eta}^{\otimes n}(\rho_x)$ and
$s=n^{-1}\sum_xp_xH(B_x)$. Conditional states have finite energy
almost surely. The thermal entropy maximum and subadditivity imply
$0\le s\le g(\eta N)$ and $H(\bar B)\le ng(\eta N)$.
Choose $\lambda\in[0,1]$ so $s=g(\eta\lambda N)$; for $N=0$
choose any $\lambda$.
Degradation with $\kappa=(1-\eta)/\eta$ and Jensen give
\begin{align}
 n^{-1}\sum_xp_xH(E_x)&\ge f_\kappa(s)
                 =g((1-\eta)\lambda N),\label{e}\\
 n^{-1}\sum_xp_xH(\rho_x)&\le
          g(g^{-1}(s)/\eta)=g(\lambda N).\label{a}
\end{align}
For \eqref{a}, average the entropy bound at $\eta$, apply convexity
of $f_\eta$, then invert. These are finite-mode statements for
arbitrary entangled inputs. Countable or measurable ensembles follow
by Jensen for integrable entropies; the finite average-energy bound
controls their integrals. At $\eta=1$, the environment is vacuum
and input and output entropies coincide.

\section{Classical, quantum and entanglement resources}
\begin{theorem}\label{cqe}
Under \eqref{energy} and the stated quantum resource conventions, the
capacity region is the closure of the union over $\lambda\in[0,1]$ of
\begin{align}
 C+2Q&\le g(\lambda N)+g(\eta N)-g((1-\eta)\lambda N),\label{c1}\\
 Q+E&\le g(\eta\lambda N)-g((1-\eta)\lambda N),\label{c2}\\
 C+Q+E&\le g(\eta N)-g((1-\eta)\lambda N).\label{c3}
\end{align}
\end{theorem}
\begin{proof}
The regularized quantum dynamic theorem \cite{QDC} and its
bosonic formulation \cite[Proposition 1]{WHG} bound the three left
sides by
\[
 \frac{H(\bar B)+\sum_xp_x[H(\rho_x)-H(E_x)]}{n},\quad
 \frac{\sum_xp_x[H(B_x)-H(E_x)]}{n},\quad
 \frac{H(\bar B)-\sum_xp_xH(E_x)}{n},
\]
with vanishing normalized coding errors. Equations \eqref{e}--\eqref{a}
give one common $\lambda$ for each block and prove the converse.
Section~\ref{approx} justifies all limiting operations.

For achievability purify the displaced thermal ensemble with
conditional mean $\lambda N$ and Gaussian displacement variance
$(1-\lambda)N$. Its average input is thermal of mean $N$.
The four entropies per use, respectively average $B$, conditional
input, conditional $B$, conditional $E$, are
$g(\eta N),g(\lambda N),g(\eta\lambda N),g((1-\eta)\lambda N)$.
The classically enhanced father protocol and noiseless conversions
\cite{CEF,QDC,WHG} attain the resulting polyhedron. Finite
bounded-cost approximations with $N'<N$, followed by closure,
enforce the exact stated budget.
\end{proof}
The joint erratum \cite{ERR} corrects generic optimization claims,
not the regularized region. No efficient optimization or generic
additivity claim is used.

\section{A zero-energy obstruction to the literal private model}
The public/private/key formula of \cite[Eq.\ (2)]{WHG} is
\begin{align}
 R+P&\le g(\eta N),\label{p1}\\
 P+S&\le g(\eta\lambda N)-g((1-\eta)\lambda N),\label{p2}\\
 R+P+S&\le g(\eta N)-g((1-\eta)\lambda N).\label{p3}
\end{align}
Its coding source has an operational inconsistency.
In \cite[Section 5, before Eq.\ (14)]{PDC}, $M,T_A$ are
generated private message and key, $K,L$ public data, $J$
consumed private communication, and $S_B$ consumed key.
Writing $U=(M,T_A,J,S_B)$ and $Z=(K,L,E^n)$, the predecoding secrecy criterion is
\begin{equation}\label{literal}
 \|\omega^{UZ}-\pi^U\otimes\sigma^Z\|_1
       \le\epsilon_{\rm sec},
\end{equation}
where $\pi$ is uniform on the whole product. It requires
consumed secrets jointly independent of public data and all
four private registers mutually independent.
A related overconstraint occurs in \cite[Eq.\ (4)]{PEPF}.
It contradicts the resource protocols stated in the sources;
the optimization erratum \cite{ERR} does not correct it.

\begin{proposition}\label{zero}
At $N=0$, \eqref{literal} and the source's reliability
condition with error $\epsilon_{\rm rel}$ imply
\[
 1\le\frac1{d_Md_{T_A}}+
            \frac{\epsilon_{\rm sec}+\epsilon_{\rm rel}}2.
\]
Thus vanishing errors exclude any nontrivial generated
private/key register, regardless of consumed resources.
The literal private region is not \eqref{p1}--\eqref{p3}.
\end{proposition}
\begin{proof}
The positive number operator has one-dimensional vacuum kernel.
Zero average energy forces every transmitted input and Bob's
channel output to be vacuum. Trace out $K,E^n$ in
\eqref{literal}. In the comparison state, uniform
$W=(M,T_A)$ is independent of Bob's full nonvacuum input
$(L,J,S_B)$. Any decoder reproduces $W$ with probability at
most $1/(d_Md_{T_A})$. Trace contraction changes that
probability by at most $\epsilon_{\rm sec}/2$.
Reliability, Eq.\ (13) of \cite{PDC}, requires probability
at least $1-\epsilon_{\rm rel}/2$.
The positive integer $d_Md_{T_A}$ must eventually be one
when both errors vanish.

Yet the displayed $N=0$ union contains one-time pad
$(-1,1,-1)$. For independent uniform $d$-symbol $M,S_B$ and
$L=M+S_B$ modulo $d$, $M$ alone is exactly secret, whereas
the minimum trace distance in \eqref{literal} is
$2(1-1/d)$. Each fixed $L$ supports only $d$ of the $d^2$
private pairs; comparison with a uniform pair gives the
distance, also minimized over the public marginal.
\end{proof}
This disproves the unchanged literal private target. It is
not an entropy-inequality failure or merely a failure of
one particular coding construction.

\section{The explicit generated-secret repair}
Keep signed forward resource accounting, the same encoding and decoding
normal form, finite gross rates and correctness. Define
$X=(K,L)$, $W=(M,T_A)$, $V=(J,S_B)$, $Y=(W,V)$ and require
\begin{equation}\label{gen}
 \|\omega^{W X E^n}-\pi^W\otimes\omega^{X E^n}\|_1\longrightarrow0.
\end{equation}
Consumed $V$ may correlate with $W$ and public transcript.
Eve does not receive the consumed key as an extra side channel.
Generated public and private communication labels are fresh uniform
messages initially independent of the preshared key. Arbitrary public
common-randomness generation is not identified with communication.
This protects uniform generated private messages and keys;
it does not claim semantic security for arbitrary message
distributions. The change from \eqref{literal} is explicit.

\begin{theorem}\label{rps}
Under \eqref{gen} and the remaining stated conventions, the
private capacity region is the closed union of
\eqref{p1}--\eqref{p3} over $\lambda\in[0,1]$.
\end{theorem}
\begin{proof}[Corrected converse]
Correctness, continuity on finite generated registers and data
processing give decoding information bounds with $o(n)$ errors.
Secrecy gives $I(W;XE^n)=o(n)$ by conditional-entropy continuity
on $W$ at fixed gross rates. It does not imply joint secrecy of $WV$.
Let overbars denote generated rates and tildes consumed rates;
$\log|L|=n\widetilde R$ and
$\log|V|=n(\widetilde P+\widetilde S)$.
The first bound is unchanged:
\begin{align*}
 n(\overline R+\overline P)
 &\le I(KM;B^nLJS_B)+o(n)\\
 &\le I(XY;B^n)+H(L)+H(J)+o(n).
\end{align*}
This uses initial independence of $KM$ from $S_B$, the chain
rule and $I(KM;LJ|S_B)\le H(LJ)$, as in the source's first bound.
Subtract the consumed public/private rates.

For the second bound, correctness and secrecy yield
$n(\overline P+\overline S)\le
I(W;B^nV|X)-I(W;E^n|X)+o(n)$.
Expand the right side as
\begin{align*}
 &I(Y;B^n|X)-I(Y;E^n|X)\\
 &\quad+I(W;V|X)-I(V;B^n|X)+I(V;E^n|WX)\\
 &\le I(Y;B^n|X)-I(Y;E^n|X)+H(V|X).
\end{align*}
Indeed the two positive remaining terms sum to
$I(V;WE^n|X)\le H(V|X)$ for classical $V$.
Subtract the consumed private/key rates.

Finally,
\begin{align*}
 n(\overline R+\overline P+\overline S)
 &\le I(KW;B^nLV)-I(W;E^n|X)+o(n)\\
 &\le I(XY;B^n)-I(Y;E^n|X)+H(L)+H(V)+o(n).
\end{align*}
Expand the first line into the desired difference,
$-I(LV;B^n)$, and
$I(KW;LV)+I(V;E^n|KWL)$.
The latter is at most
$H(L)+I(V;KWE^n|L)\le H(L)+H(V)$.
Subtract all three consumed rates.

For a degradable channel, pure spectral refinement of mixed
conditional inputs cannot decrease any private information bound:
the extra difference is
$I(Z;B^n|XY)-I(Z;E^n|XY)\ge0$ by data processing.
For pure refinements $H(B^n|XY)=H(E^n|XY)$.
The three bounds consequently reduce to at most
$H(\bar B)$, $H(B^n|X)-H(E^n|X)$, and
$H(\bar B)-H(E^n|X)$.
Equation \eqref{e} gives all three claimed photon bounds.
\end{proof}
\begin{proof}[Achievability for generated-secret security]
A finite cq ensemble $p_xp_{y|x}\rho_{xy}$ has a generated-secret
achievable corner
\[
 (R,P,S)=\bigl(I(X;B),I(Y;B|X),-I(Y;E|X)\bigr).
\]
Use HSW packing and quantum covering \cite{DEV,HLB,HN}, not the
inconsistent stronger criterion of \cite{PEPF}.
An outer public HSW code uses typical $x^n$. For each $x$ an
independent inner random table is indexed by private message $m$
and uniform consumed key $s$. Bob, given $s$, decodes below
$I(Y;B)_x$. Averaging unknown $s$ soft-covers Eve above key
rate $I(Y;E)_x$. Paste the finitely many scalar codes over
the positions of each $x$. Outer ensemble-average HSW decoding
at $I(X;B)$ and gentle measurement preserve inner decoding.
Expected error and uniform-message trace secrecy tend to zero.
This gives \eqref{gen} without any claim of secrecy after
revealing $s$. The supplement gives the precise packing,
covering and mean-cost selection.

For the bosonic ensemble take independent Gaussian
$\alpha,\beta$ of variances $(1-\lambda)N,\lambda N$ and
coherent input $|\alpha+\beta\rangle$.
The average input is thermal of mean $N$ and the conditional
input given $\alpha$ is displaced thermal of mean $\lambda N$.
Individual $B,E$ outputs are pure. Thus the corner has
$I(X;B)=g(\eta N)-g(\eta\lambda N)$,
$I(Y;B|X)=g(\eta\lambda N)$,
$I(Y;E|X)=g((1-\eta)\lambda N)$.
The three classical unit conversions translate that corner
to exactly \eqref{p1}--\eqref{p3}; generated secrets remain
protected when consumed resources are treated as consumed.
Their composition is justified by trace-norm contraction: the joint
generated secret is uniform and independent of Eve and the public
transcript, so it may supply a subsequent private-message or key
resource. Subsequent uses discard consumed registers. Finite resource
cancellation is the original net-rate convention, with any required
seed charged as a finite gross input; errors of a fixed finite
composition add and vanish. No residual secrecy for consumed key is
used.
The finite approximations and cost argument below enforce
the exact mean-energy constraint.
\end{proof}

\section{Approximation, average cost and closure}\label{approx}
For fixed finite $n$, oscillator entropy is finite and continuous
on bounded mean-energy sets; see the cutoff proof in
\cite[Lemma 2.2]{OAI}. With $P_K=\mathbf1_{\{G_n\le K\}}$, the map
$\Gamma_K(\rho)=P_K\rho P_K+
\Tr[(I-P_K)\rho]|0\rangle\langle0|$ never increases energy.
It converges in trace norm, also with references, with error at most
$2\sqrt{nN/(K+1)}+nN/(K+1)$ under \eqref{energy}.
Pure-loss output and environment of cutoff inputs remain in the
same finite number sector. In a converse, choose $K$ separately
for each block and finite resource dimension to make all decoding,
secrecy and finite-reference continuity errors arbitrarily small.
Output entropy terms retain bounded energy. Then take blocklength
limits. No uniform cutoff or input product restriction is used.

For achievability first fix $0<N'<N$. Truncate Gaussian labels
to compact sets and approximate by finite atomic distributions.
Trace continuity, bounded energy on compact sets and Gaussian
first-moment tails give convergence of average cost and all
conditional entropy quantities. Next project each conditional input
onto a one-mode number cutoff and normalize.
Its conditional mean energy cannot increase, since all removed
number levels are higher than the retained levels. States,
energies and entropies converge. Purifications compress to finite
reference dimension; projected coherent states remain pure.
For sufficiently fine approximation the mean cost $\mu<N$.

On a fixed cutoff the cost observable has norm at most $K$ per
use. Random quantum father codes have expected input trace-close
to the selected conditional product state
\cite[Propositions 2--3]{CEF}; the direct private construction
has the same property by random sampling.
Classical typical types change per-use cost by at most
$|X|\delta K$, and expected-input trace error $\epsilon_n$
changes it by at most $K\epsilon_n$.
Thus expected code cost is at most
$n(\mu+|X|\delta K+K\epsilon_n)<nN$ for small $\delta$ and
large $n$. Markov gives a fixed positive probability of cost
at most $nN$. The probability that error or secrecy exceeds
the square root of its vanishing expectation tends to zero.
Their intersection therefore contains a feasible good code.
This selects the original codes without postselecting quantum
data or changing rates. Take coding limits, approximation limits
and then $N'\uparrow N$. At $N=0$ the corrected model attains
the conversion cones with vacuum channel inputs.

Both displayed unions are closed: for convergent rate triples,
take a subsequence of their parameters in compact $[0,1]$ and
pass the inequalities by continuity. They are convex as well.
Use $s\in[0,g(\eta N)]$ instead of $\lambda$.
All right sides are concave in $s$ by Lemma~\ref{scalar},
because $g(g^{-1}(s)/\eta)$ is concave and $f_\kappa(s)$ is
convex. Their feasible sets in rate-plus-$s$ space are convex;
projection gives the unions. This handles time sharing and
all required closures.

\section{Boundaries and verification limits}
At $N=0$ every right side is zero, giving the quantum and
corrected private conversion cones. Proposition~\ref{zero}
is the different literal-model conclusion.
At $\lambda=0$ the conditional entropy terms vanish; at
$\lambda=1$ they are the full thermal values.
At $\eta=1/2$, $B,E$ have equal entropies, so $Q+E\le0$
and $P+S\le0$.
At $\eta=1$ the environment entropy is zero, giving the
energy-constrained noiseless formulas.
No argument divides by $N$, $1-\eta$ or a zero entropy.

The unchanged literal private target is refuted; the positive
private theorem has the explicit repaired convention.
There is no claim for superlinear catalysts, stronger energy
constraints, semantic security, thermal-channel quantum
capacity, amplifiers or two-way assistance.
Ordinary single-resource pure-loss capacities and ordinary
classical two-receiver broadcast capacity are not new.
The cited confidential-broadcast $\eta<1/2$ formula fails
the $\eta=0$ common-message check when Bob must decode;
no confidential-broadcast extension is included.

\paragraph{Disclosure.}
AI tools were used extensively in research, drafting and verification.
Automated independent adversarial reviews are recorded in the
supplement; they are not conventional human peer review or refereeing.
This note has not undergone conventional human peer review at
publication. The upstream EPnI is an explicitly cited unrefereed
external proved theorem whose soundness is required for its
consequences; we do not silently assume an unproved entropy premise.
No external individuals were contacted for this work.

\begin{thebibliography}{99}
\raggedright
\bibitem{OAI} OpenAI, \emph{The entropy photon-number inequality},
OpenAI Math Release preprint (2026),
\href{https://github.com/openai/math/blob/adc7f1241b42e322a6451854ab7e4b4c146bf78a/preprints/The-entropy-photon-number-inequality-September-24-2026/paper.pdf}{pinned manuscript}.
Supplied key: \texttt{OAI:The-entropy-photon-number-inequality-September-24-2026}.
\bibitem{WHG} M.~M.~Wilde, P.~Hayden and S.~Guha,
\emph{Quantum trade-off coding for bosonic communication},
Phys.\ Rev.\ A \textbf{86}, 062306 (2012),
\href{https://doi.org/10.1103/PhysRevA.86.062306}{doi:10.1103/PhysRevA.86.062306}.
\bibitem{DP} G.~De Palma, \emph{New lower bounds to the output entropy
of multi-mode quantum Gaussian channels}, IEEE Trans.\ Inf.\ Theory
\textbf{65}, 5959--5968 (2019),
\href{https://arxiv.org/abs/1805.12469v3}{arXiv:1805.12469v3}.
\bibitem{QDC} M.~M.~Wilde and M.-H.~Hsieh,
\emph{The quantum dynamic capacity formula of a quantum channel},
Quantum Inf.\ Process.\ \textbf{11}, 1431--1463 (2012),
\href{https://arxiv.org/abs/1004.0458v3}{arXiv:1004.0458v3}.
\bibitem{PDC} M.~M.~Wilde and M.-H.~Hsieh,
\emph{Public and private resource trade-offs for a quantum channel},
Quantum Inf.\ Process.\ \textbf{11}, 1465--1501 (2012),
\href{https://arxiv.org/abs/1005.3818v3}{arXiv:1005.3818v3}.
\bibitem{ERR} M.~M.~Wilde and M.-H.~Hsieh, \emph{Erratum to:
The quantum dynamic capacity formula of a quantum channel and
Public and private resource trade-offs for a quantum channel},
Quantum Inf.\ Process.\ (2012),
\href{https://doi.org/10.1007/s11128-012-0451-2}{doi:10.1007/s11128-012-0451-2}.
\bibitem{CEF} M.-H.~Hsieh and M.~M.~Wilde,
\emph{Entanglement-assisted communication of classical and quantum
information}, IEEE Trans.\ Inf.\ Theory \textbf{56}, 4682--4704 (2010),
\href{https://arxiv.org/abs/0811.4227v4}{arXiv:0811.4227v4}.
\bibitem{PEPF} M.-H.~Hsieh and M.~M.~Wilde,
\emph{Public and private communication with a
quantum channel and a secret key}, Phys.\ Rev.\ A \textbf{80},
022306 (2009), \href{https://arxiv.org/abs/0903.3920v1}{arXiv:0903.3920v1}.
\bibitem{HLB} M.-H.~Hsieh, Z.~Luo and T.~A.~Brun,
\emph{Secret-key-assisted private classical communication capacity
over quantum channels}, Phys.\ Rev.\ A \textbf{78}, 042306 (2008),
\href{https://arxiv.org/abs/0806.3525}{arXiv:0806.3525}.
\bibitem{DEV} I.~Devetak, \emph{The private classical capacity and
quantum capacity of a quantum channel}, IEEE Trans.\ Inf.\ Theory
\textbf{51}, 44--55 (2005),
\href{https://arxiv.org/abs/quant-ph/0304127}{arXiv:quant-ph/0304127}.
\bibitem{HN} M.~Hayashi and H.~Nagaoka, \emph{General formulas for
capacity of classical--quantum channels}, IEEE Trans.\ Inf.\ Theory
\textbf{49}, 1753--1768 (2003),
\href{https://arxiv.org/abs/quant-ph/0206186v4}{arXiv:quant-ph/0206186v4}, Lemma 2.
\end{thebibliography}
\end{document}
